% 主旨：结论逐条对应贡献；限定性说明集中在局限一节（单次运行、监控指标、规则与学习混合、
%       基线不完全可比、规则依赖英雄、缺少采样后改写的对照）；后续方向是多种子评测与软先验。
\section{结论与局限}
\label{sec:conclusion}

针对短训预算下平台 PPO 基线无法推塔的问题，把鲁班七号的战术规则写成采样前、训练与推理共用的合法动作掩码，并配合混合塑形奖励与冷启动偏置。主要结论如下。

\begin{enumerate}
  \item 掩码在采样前施加并随样本保存，使执行动作始终来自策略分布、新旧策略在同一支持集上计算比率；它用先验知识把策略限制在 $\Pi_{\mathcal M}$ 内，而不引入特权信息。同代码、同超参数的 \qty{30}{\minute} 冷启动训练中，完整配置对 \path{common_ai} 的终点监控胜率为 \num{1.00}；采用旧规则与权重的配置为 \num{0.96}。
  \item 推塔能力主要来自掩码：只用稀疏奖励时终点胜率为 \num{0.80}，且训练开始即为 \num{0.67}；离线固定权重、只改掩码的对比（\num{1}/\num{10} 到 \num{8}/\num{10}）也显示掩码直接改变了补给与推进行为。去掉掩码后混合奖励没有带来胜局；保留掩码时，混合奖励把胜率由 \num{0.80} 提到 \num{1.00}，己方外塔余血由 \num{3526} 提到 \num{8107}。
  \item 续训模型对 \path{common_ai} 保持全胜，伤抗比由 \num{0.95} 升至 \num{1.58}，与续训起点对战取得 \num{4} 胜 \num{2} 负。
\end{enumerate}

与 \textcite{ye2020mastering} 的大规模训练相比，本文的实验显示：当训练量缩到分钟级时，原本作为辅助手段的动作剪枝决定了策略能否推塔；把战术规则写成与策略一致的掩码，可以在不改动 PPO 损失形式的前提下直接用于训练。

\subsection{局限}
\label{sec:limitations}

第一，每组云端配置只运行了一次，组间差异中有多少来自随机种子无法估计；“完整（旧规则）”是另一套配置，不能当作“完整”的重复。离线对比的样本也只有 \numrange{5}{20} 局，六局模型对战不足以区分实力。第二，主要指标来自训练监控，对手只有内置规则智能体，监控统计按评估局滚动汇总，不等同于固定检查点上的独立评测。第三，掩码中的推进走廊直接把移动方向锁向敌塔，“仅掩码”在训练开始时就能赢下多数评估局，因此对 \path{common_ai} 的胜率衡量的是规则与学习的组合；规则之外策略学到了多少，只能从“仅掩码”与“完整”的差距间接估计。第四，“基线”使用平台官方超参数且训练 \qty{25}{\minute}，与其余各组不完全可比，缺少与本文超参数一致的“无掩码、稀疏奖励”组，因此不能确立掩码与奖励之间的一般交互关系；“续训”同时改变了规则、权重和训练量，它与起点的差异不能归因到单一改动。第五，本文没有实现采样后改写动作的对照，一致性掩码相对这种做法在训练效果上的优势只有机制上的论证。第六，掩码中的射程、血线和停靠距离按鲁班七号与 V63.0.9 任务 1 设定，换用其他英雄、加入小兵特征或改变判胜规则时需要重新标定，规则本身也限制了策略可能超越人工经验的空间。

后续可在相同预算下对每组配置运行多个随机种子，补齐“无掩码、稀疏奖励”组，并以固定检查点对多个对手做独立评测。

算力充足时，可把战术掩码换成软先验，让策略逐步走出 $\Pi_{\mathcal M}$（\cref{eq:policy-class}）；环境层面的合法性掩码仍保持硬约束，只软化战术约束。一种做法是不再把被战术规则排除的动作概率置零，而是在对数几率上加偏置：$\tilde\ell_\theta(a\mid s)=\ell_\theta(a\mid s)+\eta_t\,b(s,a)$，其中规则允许的动作 $b=0$，被排除的动作 $b=\log\xi$，$0<\xi\ll1$。$\eta_t=1$ 时被排除动作的相对概率被压低到原来的 $\xi$ 倍，接近硬掩码；$\eta_t$ 随训练衰减到 \num{0} 时先验完全撤除。另一种做法是在损失中加入被排除动作的概率质量 $\sum_{a\notin\mathcal M(s)}\pi_\theta(a\mid s)$ 作为惩罚，并逐步减小其系数。两者都在训练早期对探索分布施加偏置，后期允许策略在规则之外找到更好的动作；冷启动偏置（\cref{sec:prior}）已是这种软先验最简单的形式。检验的关键是：战术规则完全撤除后，策略能否保持推塔行为，并超过本文“完整”配置的表现。
